\documentclass[11pt]{article}
\usepackage[margin=1in]{geometry}
\usepackage{amsmath,amssymb,amsthm,booktabs}
\usepackage[hidelinks]{hyperref}
\usepackage{xurl}
\newtheorem{theorem}{Theorem}
\newtheorem{lemma}[theorem]{Lemma}
\newtheorem{corollary}[theorem]{Corollary}
\newcommand{\Z}{\mathbb Z}
\newcommand{\F}{\mathbb F}
\title{The integral Fox ideal through a fixed alternating quotient}
\author{}
\date{18 November 2026}
\begin{document}\maketitle
\begin{abstract}
The first Fox row of a four-meridian presentation of an explicit
prime knot family generates the unit ideal under a fixed quotient onto
$A_8\times C_2$. The two factors record meridian permutations and
meridian parity. We certify the two sign evaluations using conjugation
averages in $A_8$. An integral average reduces the proof to the four
primes dividing $|A_8|$. The result excludes coefficient models through
this quotient in every characteristic and extends to finite $p$-group
lifts in characteristic $p$.
\end{abstract}

Let $K_n$ be the closure of the braid
\begin{equation}\label{eq:braid}
\sigma_1\sigma_2\sigma_3^{-2}\sigma_2\sigma_3^{-2}\sigma_2
\sigma_1^2\sigma_2^3\sigma_3^{-1}\sigma_2^{3+6n},\qquad n\ge0.
\end{equation}
Write $G_n$ for its knot group. The four-strand Wirtinger presentation
uses top meridians $x_0,x_1,x_2,x_3$. The outgoing words on a positive
braid crossing are $(a,b)\mapsto(b,b^{-1}ab)$ and on a negative crossing
$(a,b)\mapsto(aba^{-1},a)$. If $F_j$ denotes the bottom word in slot
$j$, the relators are $F_jx_j^{-1}$. Each local rule preserves the
ordered product. Thus the first three closure relators imply the
fourth. Put $R_0=F_0x_0^{-1}$. The terminal $\sigma_2$ letters change
only slots $1,2$, so $R_0$ is the same free word for every $n$.
This presentation and its signed diagram binding are given in
\cite[Lemma~1]{fox}.

Use permutations of $\{0,\dots,7\}$ with multiplication
$(ab)(i)=a(b(i))$. The fixed quotient $q_n\colon G_n\to A_8$ sends the
top meridians to
\[
 p_0=(01)(23),\qquad p_1=(14)(23),\qquad p_2=(05)(16),\qquad
 p_3=(12)(67).
\]
The braid at $n=0$ returns this tuple, and the terminal six-crossing
block returns it again. The permutation calculation and the full
$A_8$ image are verified in \cite[Theorem~3]{fox}.
Let $z$ be the nonidentity element of $C_2$. Write $\widehat q_n$
for the homomorphism sending $x_i$ to $(p_i,z)$.
The $C_2$ factor is the parity of the meridian exponent sum. Enumerating
products of the four tuples $(p_i,z)$ gives $40320$ distinct elements,
so $\widehat q_n$ surjects onto $A_8\times C_2$.

The Fox ideal for the displayed presentation is the two-sided ideal
in $\Z[G_n]$ generated by the images of the Fox derivatives of its
three relators. A nontrivial matrix representation of its quotient
ring would certify ordinary rank four by the criterion reported in
\cite[Theorem~3.5]{moriah}. The calculation here restricts which
quotients can carry such a representation. All certificates below
use the four derivatives of $R_0$ alone.

\begin{theorem}\label{thm:integral}
For every $n\ge0$, the four first-row Fox images generate the whole
ring $\Z[A_8\times C_2]$ as a two-sided ideal under $\widehat q_n$.
They also generate the whole ring $\Z[A_8]$ under $q_n$. Every
nonzero ring model through either of these two maps therefore fails
to annihilate the first Fox row, over any coefficient ring.
\end{theorem}

\section{Conjugation averages and the two sign evaluations}
Write $B=\Z[A_8]$. Evaluation at $z=\epsilon$, for
$\epsilon\in\{1,-1\}$, sends the first row in $B[C_2]$ to
$d_i^\epsilon\in B$, with $0\le i\le3$. To compute these elements,
a positive letter $x_i$ contributes its prefix to the left Fox
derivative, and a negative letter $x_i^{-1}$ contributes minus the
prefix ending at that letter. Evaluate each prefix at $x_i=\epsilon p_i$.
This prescription takes place in $B$, including all integer signs.
As $R_0$ and the top images are fixed, both rows are independent of $n$.

Let $J_\epsilon=(d_0^\epsilon,\dots,d_3^\epsilon)$ be a two-sided ideal
of $B$. For $H\le A_8$ and $u=\sum_g u_g g\in J_\epsilon$, put
$T_H(u)=\sum_{h\in H}huh^{-1}$. This sum belongs to $J_\epsilon$.
For each $H$-conjugation orbit $O$ in $A_8$, the coefficient of each
member of $O$ in $T_H(u)$ is
\begin{equation}\label{eq:average}
               \frac{|H|}{|O|}\sum_{g\in O}u_g .
\end{equation}
Indeed, for any $a,b\in O$, the number of $h\in H$ conjugating $a$
to $b$ is the order of the stabilizer of $a$.
Thus the orbit totals suffice to verify a conjugation average.

\begin{lemma}\label{lem:certificates}
For each $\epsilon$, the ideal $J_\epsilon$ contains $20160\cdot1$.
Its reduction modulo each of $2,3,5,7$ contains $1$.
\end{lemma}
\begin{proof}
For each of the two signs, the finite certificate supplies an element
$u=\sum_t c_t a_t d_{i_t}^\epsilon b_t$, with integer $c_t$ and
$a_t,b_t\in A_8$. The certificate also specifies generators for $H$.
There is one certificate over the integers and one for each of
$\F_3,\F_5,\F_7$ per sign. In each case the orbit total at $\{1\}$
is $1$ and the other orbit totals vanish in the indicated ring.
The subgroup sizes, orbit counts and numbers of terms are as follows.
\begin{center}
\begin{tabular}{lrrrr}
\toprule
Ring & $|H|$ & Orbits & $\epsilon=1$ terms & $\epsilon=-1$ terms\\
\midrule
$\Z$ & 20160 & 14 & 15 & 16\\
$\F_3$ & 64 & 389 & 259 & 263\\
$\F_5$ & 1344 & 35 & 30 & 26\\
$\F_7$ & 720 & 59 & 49 & 53\\
\bottomrule
\end{tabular}
\end{center}
The integer identity is therefore $T_{A_8}(u)=20160\cdot1$.
Each prime in the other rows is coprime to the corresponding
$|H|$, so \eqref{eq:average} gives a unit scalar in the reduced ideal.

The data use literal permutation tuples for $a_t,b_t$ and the subgroup
generators. The verifier generates all subgroup elements and their
conjugation orbits, evaluates the first Fox row from the signed braid
presentation, and sums the translated coefficients in each orbit.
The exact certificate is \path{integral_a8_fox_ideal_certificates.json},
with SHA-256
\url{a4ed6fdeed2fc9262addfa08730d0c4a59a9bcfd68f0f02b29da22cd0eac770c}.
The companion verifier needs integer arithmetic and permutation
composition only.

In characteristic two the two sign evaluations coincide. The earlier
identity $X_0d_0^1+X_1d_1^1=1$ in $\F_2[A_8]$
\cite[Theorem~3]{fox} supplies the remaining case.
The verifier also replays that convolution.
\end{proof}

\begin{proof}[Proof of Theorem~\ref{thm:integral}]
By Lemma~\ref{lem:certificates}, $B/J_\epsilon$ is annihilated by
$20160=2^6\cdot3^2\cdot5\cdot7$. Since its quotient by each of
$2,3,5,7$ is zero, multiplication by each of these primes is
surjective on its underlying additive group. Their product with
multiplicities is then surjective, and is also the zero operator.
It follows that $B/J_\epsilon=0$ for both signs.

Let $J$ be the ideal of $B[C_2]$ generated by the four first-row images.
Choose $j_+,j_-\in J$ that evaluate to $1$ at $z=1,-1$ respectively.
Since the evaluation kernels are $(z-1)$ and $(z+1)$, we have
$(1-j_+)(1-j_-)=0$ in $B[C_2]$. Hence
$1=j_++j_--j_+j_-$ belongs to $J$. The identity remains true after
applying any ring homomorphism. A homomorphism that kills the four
images would therefore kill $1$ and be zero. Projection to $A_8$
recovers the first sign evaluation.
\end{proof}

\section{Finite normal \texorpdfstring{$p$}{p}-group kernels}
\begin{corollary}\label{cor:pgroup}
Suppose a finite quotient $r\colon G_n\twoheadrightarrow Q$ admits
a map onto $A_8\times C_2$ or onto $A_8$ inducing, respectively,
$\widehat q_n$ or $q_n$. If the map from $Q$ has a $p$-group kernel for a prime $p$,
then over every commutative ring $R$ in which $p$ is nilpotent, the
first Fox row of $G_n$ generates $R[Q]$ as a two-sided ideal.
\end{corollary}
\begin{proof}
First take $pR=0$, and let $N$ be the normal $p$-group kernel.
The augmentation ideal $\Delta\subset R[N]$ is nilpotent. To see this,
induct on $|N|$. For $N\ne1$, choose a central element $c$ of order $p$
using the class equation. Modulo the ideal $(c-1)$, the group algebra
is $R[N/\langle c\rangle]$. Induction gives $\Delta^a\subseteq(c-1)$
for some $a$, and the latter ideal has zero $p$th power because
$c-1$ is central and $(c-1)^p=c^p-1=0$.
The kernel of $R[Q]\to R[Q/N]$ is $R[Q]\Delta$ by collecting along
cosets of $N$. Conjugation by $Q$ preserves $\Delta$, so its powers
are $R[Q]\Delta^k$ and the kernel is nilpotent.

The ideal generated by the first row maps onto $R[Q/N]$ by
Theorem~\ref{thm:integral}. It contains a lift of $1$, which is a
unit because it differs from $1$ by an element of the nilpotent
kernel. If instead $p$ is nilpotent in $R$, first apply this argument
modulo $pR$. A lift of $1$ then differs from $1$ by an element of
$pR[Q]$, again a nilpotent ideal.
\end{proof}

These results concern the displayed presentation and quotient maps.
The ordinary rank of $G_n$ and models through other quotient kernels
remain open in this family.

\paragraph{AI assistance and review.}
ChatGPT (OpenAI) assisted with mathematical arguments, literature checks,
exposition, proof review and verification code.

\begin{thebibliography}{9}
\bibitem{fox} \emph{Uniform Fox-ideal model barriers for a prime
four-bridge family}. Companion proof, 18 November 2026.
\bibitem{moriah} Y.~Moriah, \emph{Heegaard splittings of knot
exteriors}, Geometry \& Topology Monographs \textbf{12} (2007),
191--232. \href{https://doi.org/10.2140/gtm.2007.12.191}{doi:10.2140/gtm.2007.12.191}.
\end{thebibliography}
\end{document}
